\makeabstract
    {
    Probing the Optical Effects of Azulene Intercalation on a Porous 2D Perovskite, APOSS[MnCl\textsubscript{4}]\textsubscript{4}
    }
    {
    Sequoia Post-Leon\textsuperscript{1}, Viviana Glick\textsuperscript{2}, Connor Dalton\textsuperscript{2}, Douglass Reed\textsuperscript{2}, Claire Gervais\textsuperscript{2}
    }
    {
    \textsuperscript{1} Amherst College, Amherst, MA\\
\textsuperscript{2} Department of Chemistry, University of Washington, Seattle, WA 
    }
    {
    Manganese (Mn)-based perovskites, such as porous APOSS[MnCl\textsubscript{4}]\textsubscript{4}, are excellent light emitters, despite being poor absorbers. It has been shown in other Mn-based perovskites that select organics can transfer some of their absorbed energy to the Mn lattice, allowing the perovskite to emit light without direct absorption via Forster Resonance Energy Transfer (FRET). This project examines changes in the photoluminscence emission and excitation spectra of APOSS[MnCl\textsubscript{4}]\textsubscript{4} following the intercalation of a luminescent molecule, azulene, into its pores. 
    }
    {
    APOSS[MnCl\textsubscript{4}]\textsubscript{4} was sythesized. Diffuse reflectance spectroscopy (DRS) was run to determine the primary excitation wavelengths of APOSS[MnCl\textsubscript{4}]\textsubscript{4} with and without azulene intercalated. Azulene was dissolved in isopropanol and intercalated into APOSS[MnCl\textsubscript{4}]\textsubscript{4}. The presence and approximate amount of azulene intercalated was found via nuclear magnetic resonance (NMR). Powder X-ray diffraction (PXRD) confirmed that the APOSS[MnCl\textsubscript{4}]\textsubscript{4} crystal structure was retained following azulene intercalation. Photoluminescence (PL) and photoluminescence excitation (PLE) spectroscopy was run to test azulene excitation and emission in perovskite pores and search for evidence of FRET. 
    }
    {
    PXRD confirmed that APOSS[MnCl\textsubscript{4}]\textsubscript{4} was successfully synthesized. DRS found the excitation wavelength of the Mn lattice to be 417 nm. PL at 417 nm on APOSS[MnCl\textsubscript{4}]\textsubscript{4} showed peaks at 491 nm and 605 nm, which were determined to be from APOSS-Cl\textsubscript{8} and the Mn-Cl lattice, respectively. Using varying azulene concentrations, approximately 0.5 and 1 equivalents azulene were intercalated into APOSS[MnCl\textsubscript{4}]\textsubscript{4}. Via DRS, azulene excitation was found to occur primarily at 282 nm. At 282 nm, the PL of the azulene intercalated samples showed large peaks at 376 nm, which were not present in the PL of the plain APOSS[MnCl\textsubscript{4}]\textsubscript{4}. Likewise, PLE of the azulene intercalated samples at 376 nm showed a new peak at 282 nm. PLE at 482 nm of the 0.5 equivalent azulene intercalated sample showed a novel peak at 282 nm, but no such peak appeared in the PLE at 605 nm. No peak at 605 nm was observed in the 282 nm PL in the azulene intercalated sample. 
    }
    {
    This experiment showed that a luminescent molecule can be intercalated into APOSS[MnCl\textsubscript{4}]\textsubscript{4} and that its loading amount is tuneable and concentration-dependent. Azulene was found to be able to excite and emit in APOSS[MnCl\textsubscript{4}]\textsubscript{4} pores. There was evidence of energy transfer between the APOSS-Cl\textsubscript{8} organic and azulene. There was no evidence of FRET between the azulene and Mn lattice. 
    }

    \newpage
    